\subsection{Stage 2: Retrieval-Grounded Agentic Inference and Evidence Auditing}
\label{sec:stage2}

Stage 2 operates on the quality-controlled study cohort and fixed 12-label target space established in Stage 1. Its purpose is to convert the image model's outputs into study-level decisions whose downstream modifications and supporting evidence can be reconstructed. The pipeline assigns different authority to four components: RAD-DINO supplies the primary image prediction; image-to-report retrieval supplies clinical context from visually similar training studies; a deterministic policy may modify only threshold-adjacent predictions; and two constrained large language model (LLM) calls respectively review those modifications and audit the evidence supporting the final non-absent labels. LangGraph orchestrates these components through a typed state graph. This modular design draws on retrieval-grounded and stateful LLM systems while imposing explicit decision boundaries required for clinical auditability \cite{lewis2020retrieval,guu2020realm,karpukhin2020dense,yao2023react,wu2024autogen,stateflow2024}.

The Experiment 7 execution path is
\[
\mathrm{START}\rightarrow
\mathrm{Vision}\rightarrow
\mathrm{Retrieval}\rightarrow
\mathrm{Deterministic\ Fusion+LLM\ Review}\rightarrow
\mathrm{LLM\ Evidence\ Verification}\rightarrow
\mathrm{END}.
\]

\subsubsection{RAD-DINO Study Representation and Multi-Label Prediction}

Each quality-passed view is processed using the DICOM intensity transformation and RAD-DINO input preparation defined in Stage 1. Stochastic training augmentations are disabled during inference, so a fixed DICOM study produces a deterministic model input. RAD-DINO is an image-only biomedical foundation encoder derived from the DINOv2 self-supervised representation-learning family \cite{perez-garcia_exploring_2025,oquab2024dinov2}. Report text is therefore not an input to the target study's primary prediction; it enters only after image inference as retrieved evidence. This differs from image--report representation-learning methods such as ConVIRT, GLoRIA, BioViL, and REFERS, which use paired reports as linguistic supervision during pretraining \cite{zhang2022convirt,huang2021gloria,boecking2022biovil,liu2022refers}.

For each surviving view $v$ of study $s$, the final-layer class token of the RAD-DINO Vision Transformer is used as its representation \cite{dosovitskiy2021image}:
\[
\mathbf z_{s,v}=f_{\theta}(x_{s,v})_{\mathrm{[CLS]}}
\in\mathbb R^d,
\]
where $f_{\theta}$ denotes the encoder and $d$ is its hidden dimension. The model forms one study embedding from all $n_s$ available views by concatenating element-wise mean and maximum pooling,
\[
\overline{\mathbf z}_s=\frac{1}{n_s}\sum_{v=1}^{n_s}\mathbf z_{s,v},
\qquad
\mathbf z^{\max}_{s,j}=\max_{1\leq v\leq n_s}z_{s,v,j},
\qquad
\mathbf h_s=[\overline{\mathbf z}_s;\mathbf z_s^{\max}]
\in\mathbb R^{2d}.
\]
Mean pooling retains evidence shared across projections, while maximum pooling preserves a strong feature that may occur in only one view. This choice is consistent with evidence that complementary radiographic views and study-level aggregation can improve chest-radiograph representation \cite{liu2022refers,kim2023chexfusion}. Every inferred study must contain at least one surviving view.

A dropout layer with rate $0.1$ and a linear head produce one logit for each $y\in\mathcal Y_{\mathrm{disease}}$,
\[
\mathbf a_s=W\operatorname{Dropout}(\mathbf h_s)+\mathbf b,
\qquad
p_{s,y}=\sigma(a_{s,y}),
\qquad
\sigma(u)=\frac{1}{1+e^{-u}}.
\]
The 12 sigmoid outputs are independent because multiple findings may coexist in one radiographic study \cite{wang2017chestxray8,irvin2019chexpert}. The selected checkpoint fine-tunes the final four RAD-DINO transformer blocks and freezes the preceding blocks. Let $t_{s,y}\in\{0,1\}$ be the binary target and $m_{s,y}\in\{0,1\}$ the supervision mask constructed in Stage 1, where $m_{s,y}=0$ identifies an uncertain or otherwise unsupervised study--label cell. For a mini-batch of studies $\mathcal B$, training minimizes masked binary cross-entropy,
\[
\ell_{\mathrm{BCE}}(t,p)
=-\left[t\log p+(1-t)\log(1-p)\right],
\]
\[
\mathcal L_{\mathrm{masked\text{-}BCE}}
=
\frac{
\displaystyle\sum_{s\in\mathcal B}
\sum_{y\in\mathcal Y_{\mathrm{disease}}}
m_{s,y}\,\ell_{\mathrm{BCE}}(t_{s,y},p_{s,y})
}{
\displaystyle\max\!\left(
1,
\sum_{s\in\mathcal B}
\sum_{y\in\mathcal Y_{\mathrm{disease}}}m_{s,y}
\right)
}.
\]
Thus, only cells with valid binary supervision contribute to either the numerator or the effective sample count. The $\max(1,\cdot)$ denominator reproduces the implementation's zero-valid-cell safeguard and prevents division by zero; the code evaluates the equivalent numerically stable binary-cross-entropy-with-logits form before applying the mask.

Optimization uses AdamW, which decouples parameter-wise weight decay from the adaptive gradient update rather than implementing weight decay as an $L_2$ term inside Adam's loss gradient \cite{loshchilov2019decoupled}. This separation permits the regularization coefficient to be controlled independently of the two learning rates. MEDAGENT-X assigns a smaller learning rate of $10^{-5}$ to the pretrained RAD-DINO backbone and a larger rate of $10^{-3}$ to the newly initialized classifier head, with weight decay $10^{-4}$. These numerical values are project-specific hyperparameters; they are not prescribed by the AdamW paper.

A physical batch contains two studies because each study may contribute multiple high-resolution views. Gradients are accumulated for four physical batches before each optimizer update, producing a nominal effective batch size of eight studies while remaining within accelerator memory limits. The accumulated gradient is clipped to a maximum Euclidean norm of $1.0$ before the optimizer step to limit destabilizing updates during partial backbone fine-tuning. The learning rate is increased linearly during the first $5\%$ of optimizer steps and then decays to zero along a cosine curve. Cosine annealing follows the schedule family introduced by SGDR \cite{loshchilov2017sgdr}; however, MEDAGENT-X uses a single cosine-decay trajectory and does not use SGDR's periodic warm restarts. The initial linear warm-up, its $5\%$ duration, and the absence of restarts are implementation choices rather than claims derived from that citation.

On CUDA devices, forward and backward computation use automatic mixed precision. Mixed-precision training combines lower-precision arithmetic for computationally expensive operations with higher-precision accumulation where required, reducing memory consumption and increasing throughput while retaining training stability \cite{micikevicius2018mixed}. This is particularly useful for processing multi-view studies through a large transformer encoder. Training runs for at most 15 epochs with random seed 42 and terminates early after five consecutive validation epochs without improvement. The maximum epoch count, early-stopping patience, gradient-clipping threshold, and seed are MEDAGENT-X experimental settings. The retained checkpoint is the epoch with the highest masked macro-F1 on the validation partition; no test labels participate in optimization, early stopping, threshold selection, or checkpoint selection.

\subsubsection{Label-Specific Operating Points and Gray-Zone Routing}

Because label prevalence and score distributions differ, MEDAGENT-X does not impose a universal $0.5$ decision threshold. For each label $y$, an operating threshold is selected using only supervised validation cells. Given
\[
\mathcal G=\{0.05,0.06,\ldots,0.95\},
\]
the selected threshold is
\[
\tau_y=\operatorname*{arg\,max}_{\tau\in\mathcal G}
F_1\!\left(
\{t_{s,y}:m_{s,y}=1\},
\{\mathbb I[p_{s,y}\geq\tau]:m_{s,y}=1\}
\right).
\]
The grid is traversed in ascending order and the stored value is updated only after a strict F1 improvement; ties therefore select the smallest threshold attaining the maximum. Label-specific F1 optimization is motivated by the threshold sensitivity of imbalanced classification \cite{lipton2014optimal}. The resulting image-only status is
\[
\widehat z^{\mathrm{vis}}_{s,y}=
\begin{cases}
\mathrm{present}, & p_{s,y}\geq\tau_y,\\
\mathrm{absent}, & p_{s,y}<\tau_y.
\end{cases}
\]
This procedure chooses an operating point but does not calibrate the score by temperature scaling, isotonic regression, or another post-hoc calibration method. Accordingly, $p_{s,y}$ is treated as a sigmoid model score rather than a guaranteed clinical probability \cite{guo2017calibration}.

Downstream intervention is permitted only inside a symmetric gray zone,
\[
g_{s,y}=\mathbb I\!\left[|p_{s,y}-\tau_y|\leq\delta\right],
\qquad \delta=0.15.
\]
The margin is a fixed routing policy rather than a learned reject option. It instantiates the selective-prediction principle that low-confidence cases can be routed to additional processing \cite{geifman2017selective,geifman2019selectivenet}. When $g_{s,y}=0$, the model score lies more than $\delta$ from its label-specific threshold and the prediction is treated as a strong-zone decision. Deterministic fusion must then preserve the original vision status,
\[
g_{s,y}=0
\quad\Longrightarrow\quad
\widehat z^{\mathrm{det}}_{s,y}=\widehat z^{\mathrm{vis}}_{s,y}.
\]
Because the first LLM reviews only labels changed by deterministic fusion, it cannot revise a strong-zone decision. The second LLM may still audit a final strong-zone label when that label is present, but it is non-mutating and therefore cannot alter the prediction. Thus, the gray-zone gate restricts label-changing authority to threshold-adjacent decisions while leaving stronger image-model decisions unchanged.

\subsubsection{Image-to-Report Retrieval from the Training Corpus}

MEDAGENT-X performs image-to-report case retrieval rather than text-query retrieval. General retrieval-augmented generation supplies non-parametric context to a language model \cite{lewis2020retrieval,guu2020realm,karpukhin2020dense,khandelwal2020generalization,borgeaud2022improving}; here, the query is the target study embedding and the returned text belongs to visually similar historical studies. This mechanism is related to retrieval-based chest-radiograph reporting systems \cite{endo2021retrieval,ranjit2023retrieval,jeong2024multimodal}, but MEDAGENT-X neither copies a retrieved report nor generates a report for the target study. Retrieved reports are used only for constrained label fusion and evidence attribution.

The retrieval corpus contains training studies exclusively. For every indexed training study $r$, the same fine-tuned checkpoint and mean--maximum pooling operation used at query time produce $\mathbf h_r$. Each Chroma record stores
\[
(\mathrm{study\_key}_r,\mathrm{patient\_id}_r,\mathbf h_r,D_r),
\]
where
\[
D_r=\texttt{``Findings: ''}\Vert F_r\Vert
\texttt{``Impression: ''}\Vert M_r.
\]
$F_r$ and $M_r$ are the first non-empty Findings and Impression fields encountered among the eligible rows of study $r$. Empty sections are omitted, a record with no remaining report text is rejected, and the final payload is truncated to 12,000 characters. Ground-truth label values are not placed in the vector record, ensuring that downstream reasoning receives report language rather than privileged training labels.

The persistent Chroma collection uses cosine distance \cite{chroma2024collections}. For query $s$ and candidate $r$,
\[
d_{\cos}(s,r)=1-
\frac{\mathbf h_s^\top\mathbf h_r}
{\lVert\mathbf h_s\rVert_2\lVert\mathbf h_r\rVert_2},
\qquad
\operatorname{sim}(s,r)=1-d_{\cos}(s,r).
\]
The requested neighborhood size is $K=10$. Since exclusion is performed after the vector query, the collection initially returns
\[
K'=\max(8K,K+20)=80
\]
candidates. Candidates are ordered by increasing cosine distance; the query study and all studies with the same de-identified patient identifier are then removed, and the first $K$ remaining studies form $\mathcal R_s$. Restricting the index to training studies prevents validation or test report leakage, while same-patient exclusion prevents longitudinal images from the query patient from acting as apparently independent evidence. At run initialization, the query embedding dimension is checked against the collection dimension to detect an index built with an incompatible checkpoint or pooling mode.

\subsubsection{Label-Specific Report Evidence Extraction}

Retrieved reports may contain affirmed, negated, uncertain, historical, or context-dependent statements. Prior radiology information-extraction systems therefore model negation, uncertainty, entities, and relations explicitly \cite{peng2018negbio,irvin2019chexpert,smit2020chexbert,jain2021radgraph}. MEDAGENT-X instead uses a fixed rule-based extractor at the fusion boundary so that every numerical decision input is deterministic and inspectable.

Each $D_r$ is lowercased and split on newlines, periods, semicolons, and colons. For each label $y$, a fixed dictionary $\mathcal K_y$ contains its name and selected radiological synonyms. For example, Atelectasis includes \emph{atelectasis} and \emph{volume loss}; Edema includes \emph{pulmonary edema}, \emph{interstitial edema}, \emph{vascular congestion}, and related expressions; and Lung Lesion includes \emph{lung lesion}, \emph{nodule}, and \emph{mass}. A fragment contributes at most one count to a given label, even when multiple terms in $\mathcal K_y$ match.

For a matched term beginning at character position $j$, the extractor searches the preceding 70 characters for any cue in
\[
\mathcal N=\{\text{no, without, no evidence of, negative for, absent, clear of, free of}\}.
\]
A match preceded by a cue is counted as negative; every other match is counted as positive. Let $c^+_{s,y}$ and $c^-_{s,y}$ denote the resulting counts across all $D_r$ for $r\in\mathcal R_s$. The first positive and negative fragments are retained as examples and truncated to 260 characters. This windowed rule is less linguistically expressive than dependency-based or learned report labelers \cite{peng2018negbio,smit2020chexbert}; its purpose is to expose a reproducible fusion signal. Both LLM calls receive the complete retrieved documents and may therefore identify context omitted by the count-based representation.

\subsubsection{Constrained Deterministic Gray-Zone Fusion}

Fusion is computed independently for all 12 labels. Let $a_y$ be the minimum positive count required for promotion, $b_y$ the maximum negative count allowed during promotion, and $d_y$ the minimum negative count required for demotion. The default policy is
\[
a_y=3,qquad b_y=0,qquad d_y=1.
\]
Pleural Effusion uses $a_y=4$, while Consolidation, Lung Lesion, and Lung Opacity use $d_y=3$ because their broader or overlapping terms yield noisier matches. These are project-defined policy constants, not parameters adopted from the cited report-labeling systems.

For a gray-zone vision-negative label, define promotion support as
\[
P_{s,y}=\mathbb I[c^+_{s,y}\geq a_y\land c^-_{s,y}\leq b_y].
\]
For a gray-zone vision-positive label, define demotion support and strong negative-only support as
\[
D_{s,y}=\mathbb I[c^-_{s,y}\geq d_y\land c^-_{s,y}>c^+_{s,y}],
\qquad
S^-_{s,y}=\mathbb I[c^+_{s,y}=0\land c^-_{s,y}\geq a_y].
\]
The deterministic status is
\[
\widehat z^{\mathrm{det}}_{s,y}=
\begin{cases}
\widehat z^{\mathrm{vis}}_{s,y},
& g_{s,y}=0,\\
\mathrm{present},
& g_{s,y}=1,\ \widehat z^{\mathrm{vis}}_{s,y}=\mathrm{absent},\ P_{s,y}=1,\\
\mathrm{absent},
& g_{s,y}=1,\ \widehat z^{\mathrm{vis}}_{s,y}=\mathrm{present},\ D_{s,y}=1,\ S^-_{s,y}=1,\\
\mathrm{uncertain},
& g_{s,y}=1,\ \widehat z^{\mathrm{vis}}_{s,y}=\mathrm{present},\ D_{s,y}=1,\ S^-_{s,y}=0,\\
\widehat z^{\mathrm{vis}}_{s,y},
& \text{otherwise}.
\end{cases}
\]
Consequently, a negative prediction cannot be promoted by weak or mixed report evidence; negative mentions must outnumber positive mentions before a positive prediction is softened; and full demotion to absent requires negative-only evidence meeting the promotion count bar. The pipeline records $p_{s,y}$, $\tau_y$, $g_{s,y}$, $c^+_{s,y}$, $c^-_{s,y}$, the image-only status, the deterministic status, and a human-readable rule trace for every study--label pair.

\subsubsection{First LLM: Guarded Review of Deterministic Changes}

The first LLM is a change reviewer, not an unrestricted classifier. It is called only when
\[
\mathcal C_s=
\{y:\widehat z^{\mathrm{det}}_{s,y}\neq\widehat z^{\mathrm{vis}}_{s,y}\}
\]
is non-empty, and its output must cover exactly the labels in $\mathcal C_s$. This design is related to retrieval-grounded attribution, revision, and verification \cite{gao2023rarr,dhuliawala2024chain,min2023factscore}, but the LLM is restricted to reviewing an existing deterministic intervention rather than generating a diagnosis.

For every $y\in\mathcal C_s$, the request includes the label, image-only and deterministic statuses, $p_{s,y}$, $\tau_y$, $c^+_{s,y}$, $c^-_{s,y}$, and the deterministic rule trace. It also includes every case in $\mathcal R_s$ with its rank, study key, de-identified patient identifier, distance, similarity, and complete report document. The system prompt states that these reports belong to similar patients, not to the target patient; that retrieval is indirect evidence; that no new label may be diagnosed; and that deterministic fusion should be retained unless strong label-specific contradiction is present.

Both LLM components call a locally hosted Ollama endpoint with temperature $0$ and a 300-second timeout. The model identifier remains a recorded run-time parameter. A JSON Schema is passed through Ollama's structured-output interface \cite{ollama2024structured}; constraining the generated language before programmatic use follows the same general motivation as grammar- and schema-constrained decoding \cite{scholak2021picard}. For every changed label, the first response must contain an action
\[
\mathcal A=\{\mathrm{keep},\mathrm{veto},\mathrm{uncertain}\},
\]
a corresponding final status, confidence in $\{\mathrm{low},\mathrm{moderate},\mathrm{high}\}$, an evidence assessment in $\{\mathrm{supporting},\mathrm{contradictory},\mathrm{mixed},\mathrm{insufficient}\}$, a non-empty rationale, and lists of supporting and contradicting case identifiers. Additional JSON properties are disallowed.

Local validation then enforces semantic constraints. The response may contain no missing, duplicate, or unrequested labels, and every cited identifier must belong to $\mathcal R_s$. A \emph{keep} action must preserve $\widehat z^{\mathrm{det}}_{s,y}$, a \emph{veto} must restore $\widehat z^{\mathrm{vis}}_{s,y}$, and an \emph{uncertain} action must return the uncertain status. A keep is always admissible, whereas every low-confidence non-keep action is rejected.

A veto additionally requires a contradictory assessment and at least one contradicting case. When deterministic fusion promoted an absent label to present, a veto is applied only if $c^-_{s,y}>0$, or if the LLM reports high confidence and cites at least two contradicting cases. A valid contradictory veto of a deterministic demotion may restore the original present status. An uncertainty action requires a mixed, insufficient, or contradictory assessment. After a deterministic promotion, uncertainty is applied only with high confidence and either $c^-_{s,y}>0$ or at least two contradicting cases; after a deterministic demotion, a valid uncertainty action is applied directly. Every rejected action falls back to the deterministic status and records the policy reason.

If parsing, schema validation, or guardrail validation fails, the agent issues one repair request containing the original structured request and the validation error. The repaired response is evaluated label by label: valid reviews are retained, while invalid, duplicate, or missing label reviews fall back only for their affected labels. If repair also fails, the full deterministic result is preserved. The resulting post-review status, $\widehat z^{\mathrm{fus}}_{s,y}$, is therefore never lost because an LLM is unavailable or returns malformed output. The audit record stores the deterministic and final statuses, requested action, confidence, assessment, application flag, policy or fallback reason, cited case identifiers, and rationale.

\subsubsection{Second LLM: Non-Mutating Evidence Verification}

Experiment 7 invokes a separate LLM to characterize the support for the finalized predictions. Since LLM-based evaluators can exhibit model and prompt biases \cite{zheng2023judging}, this verifier has no label-editing authority. Its role is related to work on factual and clinically significant error detection and evidence-grounded verification \cite{ostmeier2024green,gao2023rarr,min2023factscore,dhuliawala2024chain}, as well as retrieval-grounded clinical question answering \cite{zakka2024almanac}.

The verifier reviews exactly
\[
\mathcal V_s=
\{y:\widehat z^{\mathrm{fus}}_{s,y}\in
\{\mathrm{present},\mathrm{uncertain}\}\}.
\]
Thus, $\mathcal V_s$ includes unchanged strong-zone positives as well as gray-zone outputs, whereas final absent labels are excluded. For each $y\in\mathcal V_s$, the request includes $\widehat z^{\mathrm{fus}}_{s,y}$, $\widehat z^{\mathrm{vis}}_{s,y}$, $p_{s,y}$, $\tau_y$, $g_{s,y}$, whether fusion changed the label, and the complete fusion reason. The verifier receives the same ranked case records and full reports used by the first LLM.

The response schema requires, for every reviewed label, an integer evidence score from 1 to 5; confidence; separate vision-support, retrieval-support, and contradiction levels; one of the four evidence assessments; a non-empty rationale; and supporting and contradicting case identifiers. It also requires an overall study score, separate summaries of the vision, retrieval, and fusion evidence, an evidence narrative, and overall notes. Validation requires exact agreement with $\mathcal V_s$, rejects duplicates and unknown case identifiers, and requires a contradictory assessment to specify weak or strong contradiction and cite at least one contradicting case.

The verifier cannot modify $\widehat z^{\mathrm{fus}}_{s,y}$. Valid identifiers are resolved to attributable report snippets, and a parallel evidence object records the final and image-only statuses, score, support and contradiction levels, gray-zone and fusion-change indicators, rationale, and cited cases. If $\mathcal V_s=\varnothing$, the LLM is not called and the study receives an explicit no-predicted-label evidence result. If the LLM call or validation fails, Experiment 7 writes a failed-verification result with score 1 and the failure reason; it does not substitute the deterministic verifier used in Experiment 6. Unless the run explicitly permits LLM fallback, any such failure causes termination after the audit artifacts have been written.

The two calls therefore answer different questions. The first asks whether a specific deterministic change should be kept, vetoed, or softened; the second asks how well the already finalized non-absent predictions are supported. Separating decision review from evidence reporting prevents a favorable evidence narrative from becoming an implicit diagnostic action.

\subsubsection{LangGraph Orchestration, State Validation, and Provenance}

MEDAGENT-X implements the inference pipeline as a LangGraph \texttt{StateGraph}, where each node is a typed transformation over a shared study-level state \cite{langgraph2024}. Graph-based execution makes the order of operations explicit and prevents a downstream component from running before its required upstream evidence exists. This organization is related to vertex-centric graph computation, in which computation proceeds through structured state updates, and to state-machine formulations of LLM workflows \cite{malewicz2010pregel,stateflow2024}. LangGraph is used only for orchestration, state transfer, and execution tracing: the substantive decisions remain defined by the fine-tuned vision model, cosine retrieval, deterministic fusion rules, and the two guarded LLM contracts described above.

For study $s$, the graph is initialized with
\[
\mathcal S_s^{(0)}=
\{\mathrm{study\_key}_s,\mathrm{dicom\_paths}_s\}.
\]
The shared \texttt{MedAgentXInferenceState} is a typed dictionary whose optional fields are populated progressively. In addition to the two input identifiers, it can contain the structured vision result $O_s^{\mathrm{vis}}$, retrieved cases $\mathcal R_s$, the complete label-fusion-agent result $O_s^{\mathrm{agent}}$, the finalized fusion result $O_s^{\mathrm{fus}}$, the evidence-verification-agent result $O_s^{\mathrm{verifier}}$, and the final evidence record $E_s$. The distinction between $O_s^{\mathrm{agent}}$ and $O_s^{\mathrm{fus}}$ is important: the former retains both the deterministic proposal and all LLM-review metadata, whereas the latter exposes the post-guardrail labels consumed by evidence verification.

The graph is compiled once with the fixed directed path
\[
\mathrm{START}\rightarrow V\rightarrow R\rightarrow F\rightarrow E
\rightarrow\mathrm{END},
\]
where $V$, $R$, $F$, and $E$ denote the vision, retrieval, fusion-review, and evidence-verification nodes. For node $k$, let $\Delta_k(\mathcal S)$ denote the partial dictionary returned by that node. State propagation can then be written as
\[
\mathcal S_s^{(k+1)}=
\mathcal S_s^{(k)}\cup\Delta_k(\mathcal S_s^{(k)}),
\]
where newly produced fields are merged into the existing state and become available to subsequent nodes. The topology is intentionally linear because each stage has a strict evidence dependency; the graph does not permit retrieval to precede image encoding or evidence verification to precede final label fusion.

The vision node $V$ requires \texttt{study\_key} and \texttt{dicom\_paths}. It invokes the configured RAD-DINO backend and adds $O_s^{\mathrm{vis}}$, which contains the study embedding, label scores, thresholds, image-only statuses, model identifier, patient identifier, split, and source paths. Before returning, the node verifies that the study key in the model output equals the requested key. This prevents a cached or asynchronously produced prediction from being attached to the wrong patient study.

The retrieval node $R$ requires $O_s^{\mathrm{vis}}$ and uses only its study embedding and source identifiers to query the training-only Chroma collection. It adds the ranked set $\mathcal R_s$, including each neighbor's study and patient identifiers, distance, similarity, and report document. Before cohort execution, the query-embedding dimension is compared with the stored collection dimension. A mismatch causes execution to stop rather than querying an index constructed from an incompatible checkpoint or pooling representation.

The fusion node $F$ requires both $O_s^{\mathrm{vis}}$ and $\mathcal R_s$. It first converts the structured vision output into the 12 label-specific inputs required by deterministic fusion, computes $\widehat z^{\mathrm{det}}_{s,y}$, and conditionally invokes the first LLM for the changed-label set $\mathcal C_s$. It adds both $O_s^{\mathrm{agent}}$ and $O_s^{\mathrm{fus}}$ to the state. Consequently, downstream nodes receive the finalized post-guardrail statuses without discarding the deterministic baseline, requested LLM actions, validation outcomes, or fallback reasons that produced them.

The evidence node $E$ requires $O_s^{\mathrm{vis}}$, $\mathcal R_s$, and $O_s^{\mathrm{fus}}$. It invokes the second LLM only for the final non-absent label set $\mathcal V_s$ and adds both the complete verifier-agent result and the non-mutating evidence record $E_s$. The finalized labels are read-only at this boundary. Failures in either LLM stage are handled by the local policies defined above and are represented explicitly in the state through request, success, application, and fallback fields rather than being silently omitted.

Each node validates the presence of its required keys before performing computation. A missing dependency therefore produces an immediate error at the responsible graph boundary instead of propagating an incomplete object into a later prompt or decision rule. Content-level checks complement these structural checks: study identifiers must agree across the initial state and vision output, retrieved case identifiers cited by an LLM must belong to $\mathcal R_s$, response label sets must match their requested sets exactly, and all proposed label transitions must satisfy the deterministic guardrails. The graph therefore coordinates the agents, while typed contracts and local validators determine whether each state transition is admissible.

The compiled graph is invoked independently for every study. After an invocation, the runner writes one JSON Lines trajectory record and flushes it immediately. Each trajectory summarizes the study identifier, elapsed time, number of retrieved cases, deterministic and final changed-label sets, final predicted labels, LLM request and success indicators, fallback indicators, and the overall evidence score. Immediate flushing limits the amount of provenance lost if a long cohort run is interrupted.

More detailed artifacts are exported in parallel: study-level vision predictions preserve scores, thresholds, statuses, model identifiers, and source paths; label-level fusion tables preserve vision, deterministic, and finalized statuses together with rule traces; retrieval tables preserve neighbor ranks, similarities, identifiers, and report documents; and evidence-verification JSON and CSV files preserve per-label and study-level audit judgments. Separate policy-audit summaries aggregate LLM actions, validation failures, and fallback behavior. The run configuration records the checkpoint and threshold policy, vector-index path and collection name, evaluated split, hardware and mixed-precision settings, $K$, $\delta$, graph-node sequence, LLM model and endpoint, temperature, timeout, policy-version identifiers, request and success counts, fallback counts, and total elapsed time.

This provenance design provides two levels of reconstruction. At the study level, the trajectory and linked output artifacts recover the sequence of model prediction, retrieval, fusion review, and evidence verification. At the label level, the stored score, threshold, gray-zone indicator, retrieval counts, deterministic rule trace, LLM action, cited cases, validation result, and final status explain why $\widehat z^{\mathrm{fus}}_{s,y}$ was retained or changed. The combination of typed in-memory state, persisted source identifiers, explicit transition rules, schema-validated LLM outputs, and versioned run configuration therefore makes the final decisions auditable without treating LangGraph itself as a diagnostic or reasoning model.
